\chapter{Лабораторная работа №5. Конфигурирование основных сетевых служб.}
\label{ch:chapter 5}

\section{\textit{Лабораторная работа №5.1. Настройка FTP.}}

\subsection{Цели}

Лабораторная работа помогает получить практические навыки по изучению
следующих тем:

\begin{itemize}

\item Установление FTP-соединения

\item Настройка параметров FTP-сервера

\item Процедура передачи файлов на FTP-сервер

\end{itemize}

\subsection{Топология сети}

Необходимо на R1 выполнить операции с конфигурационным файлом R2.
R1 функционирует как FTP-клиент, а R2 — как FTP-сервер.

\includegraphics[width=0.5\textwidth]{lb5_1.png} 
\captionof{figure}{Топология сети для лабораторной работы №5.1}

\subsection{План работы}

\begin{itemize}

\item 1. Настройка функции и параметров FTP-сервера.

\item 2. Настройка локальных пользователей FTP.

\item 3. Вход в систему FTP-сервера с FTP-клиента.

\item 4. Выполнение операций с файлами в FTP-клиенте.

\end{itemize}

\subsection{Процедура конфигурирования}

\textbf{Шаг 1.}Настроим основные параметры устройств.

Настройте IP-адреса устройств:

\includegraphics[width=0.5\textwidth]{lb5_2.png} 

\includegraphics[width=0.5\textwidth]{lb5_3.png} 

Сохраним конфигурационный файл для последующей проверки:

\includegraphics[width=0.5\textwidth]{lb5_4.png} 

\includegraphics[width=0.5\textwidth]{lb5_5.png} 

Выведем на экран текущий список файлов:

\includegraphics[width=0.5\textwidth]{lb5_6.png}

\includegraphics[width=0.5\textwidth]{lb5_7.png}

\textbf{Шаг 2.}Настроим функцию и параметры FTP-сервера на R2.

\includegraphics[width=0.5\textwidth]{lb5_8.png}

\textbf{Шаг 3.}Настроим локальных пользователей FTP.

\includegraphics[width=0.5\textwidth]{lb5_9.png}

\includegraphics[width=0.5\textwidth]{lb5_10.png}

\textbf{Шаг 4.} Выполним вход в систему FTP-сервера с FTP-клиента.

Выполним вход в FTP-клиент:

\includegraphics[width=0.5\textwidth]{lb5_11.png}

\textbf{Шаг 5.} Выполним файловые операции в файловой системе на R2.

Настроим режим передачи:

\includegraphics[width=0.5\textwidth]{lb5_12.png}

Загрузим конфигурационный файл:

\includegraphics[width=0.5\textwidth]{lb5_13.png}

Удалим конфигурационный файл:

\includegraphics[width=0.5\textwidth]{lb5_14.png}

Выгрузим конфигурационный файл:

\includegraphics[width=0.5\textwidth]{lb5_15.png}

Закроем FTP-соединение:

\includegraphics[width=0.5\textwidth]{lb5_16.png}

\textbf{Вопросы.} 1. В активном или пассивном режиме работает FTP по умолчанию?

-В активном.




\section{\textit{Лабораторная работа №5.2. Конфигурирование DHCP.}}

\subsection{Цели}

Лабораторная работа помогает получить практические навыки по изучению
следующих тем:

\begin{itemize}

\item Настройка пула адресов на DHCP-сервере

\item Настройка глобального пула адресов на DHCP-сервере

\item Использование DHCP для статического назначения IP-адресов

\end{itemize}

\subsection{Топология сети}

Необходимо на R1 выполнить операции с конфигурационным файлом R2.
R1 функционирует как FTP-клиент, а R2 — как FTP-сервер.

\includegraphics[width=0.5\textwidth]{lb5_17.png} 
\captionof{figure}{Топология сети для лабораторной работы №5.2}

\subsection{План работы}

\begin{itemize}

\item 1. Настройка DHCP-сервера.

\item 2. Настройка DHCP-клиентов.

\end{itemize}

\subsection{Процедура конфигурирования}

\textbf{Шаг 1.} Настроим основные параметры.

Настроим их на R2:

\includegraphics[width=0.5\textwidth]{lb5_18.png} 

\textbf{Шаг 2.} Включим функцию DHCP.

\includegraphics[width=0.5\textwidth]{lb5_19.png} 

\includegraphics[width=0.5\textwidth]{lb5_20.png} 

\includegraphics[width=0.5\textwidth]{lb5_21.png} 

\textbf{Шаг 3.} Настройте пул адресов.

\includegraphics[width=0.5\textwidth]{lb5_22.png} 

\textbf{Шаг 4.} Включим функцию DHCP-сервера на GE0/0/3 R2 для назначения IP-адреса R3.

\includegraphics[width=0.5\textwidth]{lb5_23.png} 

\textbf{Шаг 5.} Настроим DHCP-клиенты.

\includegraphics[width=0.5\textwidth]{lb5_24.png} 

\includegraphics[width=0.5\textwidth]{lb5_25.png} 

Проверим настройку:

\includegraphics[width=0.5\textwidth]{lb5_26.png} 

\includegraphics[width=0.5\textwidth]{lb5_27.png} 

\includegraphics[width=0.5\textwidth]{lb5_28.png} 
